\documentclass[12pt,a4paper]{article}
\usepackage[margin=2.3cm]{geometry}
\usepackage{polyglossia}
\setmainlanguage[numerals=western]{arabic}
\setotherlanguage{english}
\newfontfamily\arabicfont[Script=Arabic,Scale=1.08]{Amiri}
\newfontfamily\englishfont[Scale=0.92]{DejaVu Serif}
\usepackage{amsmath,amssymb}
\usepackage{booktabs}
\usepackage{enumitem}
\usepackage{hyperref}
\hypersetup{colorlinks=true,linkcolor=black,urlcolor=blue,
            pdftitle={ملخص البحث},pdflang={ar}}

\newcommand{\lr}[1]{\textenglish{#1}}
% bidi reorders multi-digit runs, so section 10 printed as "01";
% force an LTR embedding on every auto-generated number.
\renewcommand{\thesection}{\LR{\arabic{section}}}
\renewcommand{\thesubsection}{\LR{\arabic{section}.\arabic{subsection}}}
\renewcommand{\thetable}{\LR{\arabic{table}}}
\newcommand{\D}{\lr{$D$}}
\newcommand{\E}{\lr{$E(P_D)$}}

\title{\bfseries طريقة هجينة لتوليد المجموعة الكفؤة الكاملة\\[2pt]
\large لبرنامج خطّي صحيح متعدّد المعايير،\\
مع معيار قرارٍ خطّي كسري}
\author{}
\date{}

\begin{document}
\maketitle

\begin{abstract}
\noindent
هذا ملخَّصٌ عربيٌّ لكلّ ما تمّ الوصول إليه: المسألة، والفكرة التي تحلّها،
والنتائج النظرية الثلاث، والقياسات العددية بحسناتها وسيّئاتها، والنتائج
السلبية، وتدقيق الصحّة الذي أُجري على البرنامج نفسه. الوثائق التفصيلية
بالإنجليزية في المستودع، وهذا الملف يحيل إليها.
\end{abstract}

\section{المسألة}

لتكن \lr{$D = \{x \in \mathbb{Z}^n_+ : Ax \leqslant b\}$} مجموعةً مقبولةً غير
فارغة ومحدودة، ولتكن المعايير المطلوب تعظيمها نِسَباً من دوالّ تآلفية:
\begin{equation*}
(P_D)\quad \text{``}\max\text{''}\ Z_k(x)
= \frac{c_k^{\top}x + a_k}{d_k^{\top}x + b_k},\quad k = 1,\dots,p,
\qquad x \in D
\end{equation*}
بشرط أن يبقى كلّ مقام \lr{$D_k(x) = d_k^{\top}x + b_k$} موجباً تماماً على \D.
والنقطة \lr{$x \in D$} \textbf{كفؤة} إذا لم توجد \lr{$y \in D$} تحقّق
\lr{$Z(y) \geqslant Z(x)$} مع متباينةٍ صارمةٍ واحدة على الأقل. ونرمز للمجموعة
الكفؤة بـ\E.

والمسألة الثانية، وهي الأصعب: بإعطاء معيار قرارٍ خطّي كسري
\lr{$\Phi(x) = (U^{\top}x + \alpha)/(V^{\top}x + \beta)$}،
\begin{equation*}
(P_E)\qquad \max\ \Phi(x) \qquad \text{\textarabic{بحيث}} \qquad x \in E(P_D)
\end{equation*}
والصعوبة بنيوية: \E\ اتحادُ أوجهٍ من \D، غير محدَّبة، وليس لها وصفٌ صريح
بمتباينات، فلا يمكن تمريرها إلى حلّال.

\paragraph{ملاحظة تحكم كل ما يلي.}
\lr{$\Phi$} دالةٌ في \lr{$x$} و\textbf{ليست} دالةً في \lr{$Z(x)$}. فنقطتان
كفؤتان لهما المتجه المعياري نفسه قد تختلف قيمة \lr{$\Phi$} عندهما. وهذا
بالتحديد سبب أن نقل البحث في فضاء المعايير من الدوال الخطية إلى الكسرية ليس
نقلاً مباشراً.

\section{الفكرة: فضاء المعايير بدل فضاء القرار}

الطرق الموجودة تعمل في \textbf{فضاء القرار}: تقتطع \D\ بقطوع
\lr{Sylva--Crema}. وكلّ قطعٍ من هذه فصلٌ منطقي \lr{(disjunction)}، فيكلّف
\lr{$p$} متغيّراً ثنائياً و\lr{$p+1$} سطراً من نوع \lr{big-$M$}. والنتيجة أن
\textbf{كل مسألة فرعية أكبر من سابقتها}، والنموذج ينمو بلا توقّف.

وما نفعله هو الاحتفاظ بالجزء غير المستكشَف في \textbf{فضاء المعايير}، كقائمة
صناديق. فالاقتطاع يصبح طَرْحَ صندوقٍ من صندوق، وكل مسألة فرعية هي \D\ مضافاً
إليها حفنةُ أسطرٍ خطّية عادية، و\textbf{النموذج لا ينمو أبداً}.

وتُبنى هذه على كمّية واحدة. بعد تحجيم معطيات كل معيار إلى أعداد صحيحة (بضرب
البسط والمقام بالعامل نفسه، فلا تتغيّر النسبة) نضع
\begin{equation*}
e_k(x;\bar x) = D_k(\bar x)\,D_k(x)\,
\bigl(Z_k(x) - Z_k(\bar x)\bigr)
\end{equation*}
فتكون \textbf{خطّيةً في \lr{$x$} وصحيحةَ القيمة}، ومن ثَمّ
\lr{$Z_k(x) > Z_k(\bar x) \iff e_k(x;\bar x) \geqslant 1$} بالضبط. وهذه هي
النقطة التي تجعل الانتقال من المعايير الخطية إلى الكسرية \textbf{أميناً لا
تقريبياً}: لا حاجة إلى تقدير أصغر خطوة.

\section{النتائج النظرية الثلاث}

\subsection{نظرية الشريحة: المجموعة الكفؤة الكاملة رخيصة}

الطريقة البديهية لإكمال الجبهة إلى مجموعة النقاط كاملةً هي فحص كفاءة كل نقطة
في كل شريحة. وهذا كلّه غير لازم:

\begin{quote}
\textbf{إذا كان المتجه \lr{$v$} غير مُهيمَنٍ عليه، فإنّ \emph{كلّ}
\lr{$x \in D$} يحقّق \lr{$Z(x) = v$} نقطةٌ كفؤة.}
\end{quote}

\noindent
والبرهان سطران: لو هيمنت \lr{$y \in D$} على مثل هذه الـ\lr{$x$} لكان
\lr{$Z(y) \geqslant Z(x) = v$} مع صرامةٍ واحدة، أي أنّ \lr{$v$} مُهيمَنٌ عليه،
وهذا خلاف الفرض.

ويتبع ذلك أمران هما مجمل المكسب: المجموعة الكفؤة هي بالضبط اتحاد شرائح
الجبهة،
\begin{equation*}
E(P_D) = \bigl\{\, x \in D \;:\; Z(x) = v
\ \ \text{\textarabic{لمتجهٍ غير مُهيمَنٍ عليه}}\ \ v \,\bigr\}
\end{equation*}
و\textbf{لا يُدفَع ولا فحصُ كفاءةٍ واحدٍ على أيّ نقطةٍ من نقاط الشرائح}.

\subsection{الشريحة خطّية، حتى للمعايير الكسرية}

الشريحة هي \D\ مضافاً إليها \lr{$p$} \textbf{معادلة خطّية}، والمفاجأة أن هذا
يصحّ للمعايير الكسرية كما يصحّ للخطّية: المعادلة \lr{$Z_k(x) = v_k$} تُكتب
\begin{equation*}
(c_k - v_k d_k)^{\top} x = v_k b_k - a_k
\end{equation*}
وهي خطّية في \lr{$x$}. وحلّ هذه المعادلات بحذف \lr{Gauss} في الحساب الكسري
المضبوط يُخرج \lr{$r$} متغيّراً \emph{قبل} أن يبدأ المسح، فينزل المسح من
\lr{$(ub+1)^n$} إلى \lr{$(ub+1)^{n-r}$}.

\textbf{القياس:} هذا هو أكبر مكسبٍ منفردٍ في العمل كلّه ---
\lr{$49.2\times$} وسيطاً على \lr{$14$} نموذجاً، وحتى \lr{$129.9\times$}،
بمجموعات نقاطٍ متطابقةٍ في الحالتين.

\subsection{الهجين الذي يدفع}

مشيٌ محليٌّ من نوع \lr{Pareto} مع أرشيف يُوَلِّد نقاطاً كفؤةً
\textbf{مُتحقَّقاً منها}، فتُوضَع في قائمة الصناديق \textbf{دون سَبْر}. وهنا
موضع الصحّة: الأرشيف مُرشَّحٌ بالهيمنة \emph{بين ما رآه فقط}، وهذا ليس فحص
الكفاءة المضبوط؛ فلا يَعبُر شيءٌ دون تحقّق. والمشي يقرّر \emph{أين ننظر}، ولا
يقرّر أبداً ما هو صحيح. والفجوة ليست نظرية: واحدٌ من كل \lr{$2000$} عضوٍ في
الأرشيف تبيّن أنه غير كفؤ، على نحوٍ قابلٍ للتكرار.

\paragraph{والترتيب هو الفرق كلّه.}
هذه أهم نتيجةٍ عمليةٍ في المشروع: الهجين يدفع \textbf{فقط} إذا سُلِّمت البذور
بالترتيب الذي كان سَبْرُ الطريقة المضبوطة سيلقاها به. أمّا بالترتيب الذي
يُنتجها الأرشيف فالهجين \textbf{أبطأ}. السطر نفسه، والنتيجة معكوسة.

\section{النتائج العددية}

\subsection{الهامش الأساسي}

\paragraph{ما هو خط الأساس، وما قيمة الهامش تبعاً لذلك.}
الجانب القراري في كل مقارنةٍ هنا هو \textbf{تنفيذنا نحن} للطريقة المنشورة، لا
كودُ أصحابها، والنماذج نماذجُنا لا نماذجَهم المرجعية. ويجب قول هذا صراحةً لأنه
يغيّر معنى الأرقام.

الجانبان يتقاسمان \lr{simplex} واحداً، وحساباً كسرياً مضبوطاً واحداً، وتفريعاً
وتحديداً واحداً، وآلةً واحدةً، ونماذجَ واحدة؛ والتنفيذ يتبع الخوارزمية
المطبوعة خطوةً بخطوة، مع تصحيح الموضعين اللذين يتناقض فيهما النصّ المطبوع مع
نفسه، والتصحيحان موثّقان. فالمقارنة إذن \textbf{تعزل الفرق الخوارزمي}، وهو
سؤالٌ أنظفُ من مقارنة كودٍ بكود --- إذ تلك تقيس أيضاً لغةَ البرمجة وخلفيةَ
الحلّال وجهدَ الهندسة.

وثمن هذا الاختيار أنّ \textbf{الثواني المطلقة لا تحمل أيّ معلومةٍ خارج هذا
التنفيذ}: \lr{Python} خالصٌ فوق كسورٍ مضبوطة أبطأُ بكثيرٍ من أيّ حلّالٍ
إنتاجي، على الجانبين سواءً بسواء. فالنسب وحدها تنتقل، والنسب وحدها تُذكَر.

\textbf{وما لم يُنجَز ولا يُدَّعى}: مقارنةٌ مع تنفيذات أصحاب الأوراق أنفسهم على
نماذجهم المرجعية. وإلى أن تُنجَز، القراءة الصحيحة لكل هامشٍ هنا هي «هذه
الخوارزمية تتفوّق على تلك حين تُنفَّذان بالطريقة نفسها»، لا «هذا العمل يتفوّق
على تلك الأوراق».


\begin{table}[h]
\centering
\caption{الهجين مقابل الطريقة المضبوطة الخالصة. أجوبةٌ متطابقة في كل الحالات.}
\begin{tabular}{lcl}
\toprule
\textarabic{التجربة} & \textarabic{النماذج} & \textarabic{النسبة} \\
\midrule
\textarabic{\E\ كاملة، \lr{$n = 6\dots12$}} & \lr{$120$}
  & \textarabic{\lr{$\mathbf{1.36\times}$}، أسرع في \lr{$116$}} \\
\textarabic{معايير كسرية فعلاً} & \lr{$30$} & \lr{$1.39\times$} \\
\textarabic{خمس عائلات بنيوية} & \lr{$150$}
  & \textarabic{\lr{$1.23\times$} وسيطاً، أسرع في \lr{$136$}} \\
\textarabic{فضاء القرار إلى فضاء المعايير} & ---
  & \textarabic{\lr{$1.5\times$} عند \lr{$p=2$}، \lr{$14.6\times$} عند \lr{$p=6$}} \\
\bottomrule
\end{tabular}
\end{table}

\noindent
الرقمان \lr{$1.36\times$} و\lr{$1.23\times$} من \textbf{حملتين مختلفتين} لا
من حملةٍ واحدة، ولا يُجمعان: الأولى على \lr{$120$} نموذجاً بتحجيم المشي بحسب
النموذج، والثانية على \lr{$150$} نموذجاً موزّعةً على خمس عائلات. والأهم أن
\textbf{كل عائلةٍ من الخمس تدفع}، بنسبٍ بين \lr{$1.14\times$}
و\lr{$1.37\times$}، فالتفوّق ليس أثراً جانبياً لعائلةٍ واحدة.

\subsection{حيث يفشل، ولماذا}

انهيار التغطية حقيقيٌّ وشديد حيث تكون الجبهة كبيرة: عند جبهةٍ من
\lr{$|F| = 107$} تنزل النسبة إلى \lr{$1.14\times$}. والسبب ليس أنّ عدد
السَّبْرات يتجاوز حجم الجبهة --- هذه فرضيةٌ اقترحتُها والقياس دحضها، إذ بقي عدد
السَّبْرات لكل متجهٍ ثابتاً بين \lr{$2.4$} و\lr{$5.2$} --- بل هو انهيار
التغطية نفسها. وتحجيم المشي بحسب النموذج يعالج ذلك: من \lr{$1.00$} إلى
\lr{$1.53\times$} عند \lr{$n=12$}.

\subsection{قابلية التوسّع}

دراسةٌ مستقلّة تُبرهن الأمثلية دون بناء \D\ ولا \E، على نماذج من
\lr{$n=10$} إلى \lr{$n=20$}: كلّها مُبرهَنة. والصعوبة تتبع \lr{$|D|$} وعدد
دورات القطع لا \lr{$n$} وحده --- النموذج نفسه عند \lr{$n=16$} يُحَلّ في ثوانٍ
إن كان المجال ضيّقاً، ويبقى بعيد المنال إن كان واسعاً.

\section{النتائج السلبية، وهي نصف العمل}

جُرِّب \lr{$22$} تغييراً: \lr{$12$} ربحت، و\lr{$9$} خسرت، وواحدٌ تعادل. ونحن
نذكر الخسارات بالطول نفسه الذي نذكر به المكاسب لأنها تحمل المعلومة نفسها.
وأبرزها:

\begin{itemize}[leftmargin=1.6em,itemsep=2pt]
\item \textbf{أربعة تغييرات تحذف ربع المسائل الفرعية ولا تشتري شيئاً.} حذف
      العمل لا يمكن أن يُكلِّف وقتاً، فالأرضية الضوضائية للآلة
      (\lr{$\pm 30\%$} لكل نموذج) لا تستطيع حلّ أثرٍ بـ\lr{$6\%$}، ولا يُصلح
      ذلك أيُّ عددٍ من التكرارات.
\item \textbf{شهادة تمامٍ ببرنامجٍ واحد:} سليمةٌ تماماً، و\lr{$4.5\times$}
      أبطأ.
\item \textbf{بذرة \lr{Tchebychev}:} غسلٌ لا يدفع أبداً (\lr{$1.02\times$}) ---
      لكنها \textbf{كمُوَلِّدٍ} أقوى من أي مجموعٍ موزون، إذ تصل إلى النقاط
      الكفؤة غير المدعومة.
\item \textbf{التسليم المضبوط--المضبوط:} تعادلٌ عند بذرةٍ واحدة، وخسارةٌ عند
      ثلاث.
\item \textbf{فكرتان تُغذّيان الآليةَ نفسها لا تتجمّعان.} الخروج المبكّر يوفّر
      \lr{$25\%$} من الصناديق وحده، و\lr{$4\%$} فقط فوق البذرة: فهو بديلٌ عنها
      لا إضافةٌ إليها. وهذا أقبل النتائج للتعميم.
\end{itemize}

\paragraph{ادّعاءٌ سُحِب.}
كتبتُ في أربع وثائق أن هناك \emph{سقفاً صلباً} للمكسب عند \lr{$1.4\times$}
تقريباً، مستنبطاً من حساب السَّبْرات. وعند \lr{$p=6$} تجاوز الزمن المقيس
\lr{($1.24\times$)} ما يسمح به ذلك السقف \lr{($1.15\times$)}. فالادّعاء
\textbf{مسحوب} وموسومٌ بأنه مسحوب في كل وثيقة؛ وحدُّ \emph{عدد} السَّبْرات
وحده هو ما بقي صحيحاً.

\section{تدقيق الصحّة: الطريقة ليست برهانها}

دقّقنا البرنامج في مقابل فرضياته المُعلَنة، فوجدنا \textbf{أحد عشر} موضعاً
فرضيةً مُعلَنةً في النصّ لم تُفحَص قطّ في الكود. أُصلحت الأحد عشر كلّها، ولها
اختبارات انحدار، ولم يتحرّك أيُّ رقمٍ مقيس. والنمط أهمّ من أيٍّ منها
منفرداً، وهو نصفان:

\paragraph{الأول: نتيجةُ برهانٍ محسوبةٌ من مسار التنفيذ.}
كل علمٍ \lr{(flag)} يدّعي أنّ شيئاً \emph{بُرهِن} كان محسوباً من شكل الحلقة ---
انتهت، ولم تبق مسألةٌ فرعية --- لا من فرضيات البرهان. فـ«قائمة الصناديق فرغت»
كانت تُبلَّغ على أنها «الجبهة تامة»، وهذا سليمٌ بالضبط إذا كانت \D\ محدودةً
والمقامات موجبة، وهذان كانا مفترَضين لا مُثبَتين. والفشل صامتٌ في الحالتين:
على \D\ غير محدودة لكنها مقبولة وذات نقاطٍ كفؤة، أعادت الطريقة
\textbf{جبهةً فارغةً موسومةً بأنها تامة}؛ وعلى معيارٍ يتغيّر إشارةُ مقامه،
أعادت متجهاً واحداً كالجبهة كلها بينما تهيمن عليه نقطةٌ مقبولة.

\paragraph{الثاني، وهو القابل للتعميم: كل توفيرٍ كان افتراضَ ثقة.}
كلّ مُعامَلٍ أُضيف إلى الواجهة البرمجية لتجنّب إعادة الحساب --- صندوق متغيّرات
محسوبٌ سلفاً، جبهةٌ محسوبةٌ سلفاً، مجموعة بذور، حلٌّ حاليٌّ مع قيمته --- هو
طريقةٌ ليُمِدّ المُنادي فرضيةً دون أن يبرهنها. أُضيف كلٌّ منها كتحسينٍ للأداء،
وصار كلٌّ منها بهدوءٍ افتراضَ ثقة. والصيغة العملية للدرس:

\begin{quote}
\textbf{كلّ مُعامَلٍ يُجنِّب المُنادي عملاً كان المُنادى سيؤدّيه، يُجنِّبه
برهاناً كان المُنادى سيملكه. وهما المُعامَل نفسه --- وأضفنا خمسةً منها دون أن
ننتبه.}
\end{quote}

\paragraph{ولماذا لم تَرَ الاختباراتُ شيئاً من ذلك.}
الاختبارات الثلاثة والتسعون كلّها كانت تُغذّي مُدخَلاتٍ \textbf{صحيحة}. فهي
تُثبت أن الطريقة تُعطي الجواب الصحيح حين تصحّ فرضياتها --- وهذا لم يكن موضع
سؤال --- ولا تستطيع أن تُثبت أنّ شيئاً يُلزِمها. ومجموعةٌ بهذا الشكل عمياءُ
عن هذه الفئة كلّها بالبناء، والعمى غير مرئيٍّ من داخلها: كل اختبارٍ ينجح، في
كل إيداع، لسنوات. وصارت المجموعة \lr{$103$} اختباراً، والعشرة المضافة كلّها
تُغذّي مُدخَلاتٍ تخرق فرضيةً وتشترط \textbf{الرفض} لا الجواب.

\subsection{أخطر اثنتين: في أدوات التحقّق نفسها}

\textbf{العاشرة} في شهادة التمام، التي تبحث عن المنافسين داخل صندوقٍ
يُسلَّم إليها. والصندوق الأصغر من \D\ يُخفيهم، فتعود الشهادة
\emph{صحيحة} على نقطةٍ ليست مثلى: بالصندوق \lr{$[1,1]$} صدّقت نقطةً كفؤةً
عند \lr{$\Phi = 1/9$} والقيمة المثلى الحقيقية \lr{$5/17$}. وهذا أسوأ موضعٍ
ممكن، لأن الشهادة هي ما يتحقّق من الأجوبة بعد الحجم الذي تعجز عنه القوة
الغاشمة --- فالبرهان الكاذب ليس جواباً خاطئاً بل \textbf{فحصاً خاطئاً على
أجوبةٍ أخرى}.

\textbf{والحادية عشرة} هي الخلل نفسه درجةً أبعد: في طريقتَي المرجع نفسيهما،
المُناداتَين من نحو \lr{$45$} موضعاً، كلٌّ يمرّر صندوقاً مكتوباً باليد. وعند
تمرير صندوقٍ أصغر من \D\ لم تفشلا ولم تُحذِّرا، بل أجابتا بثقةٍ عن المجموعة
الكفؤة لمجموعةٍ مقبولةٍ \emph{أخرى}، تحت اسمَي \E\ و«القيمة المثلى المضبوطة».
على \lr{$x_1 + x_2 \leqslant 6$} و\lr{$x_1 - x_2 \leqslant 2$}، التي مجموعتها
الكفؤة خمس نقاطٍ وقيمتها المثلى \lr{$2/5$}، أعطى الصندوق \lr{$[1,1]$} نقطةً
واحدةً سُمّيت \E، و\lr{$1/3$} سُمّيت القيمة المثلى.

\begin{quote}
\textbf{المرجع المقصوص لا يُعطي جواباً خاطئاً، بل يُعطي معياراً خاطئاً ---
والطريقة المقيسة عليه تبدو صحيحةً بالضبط حين تتّفق مع القصّ.}
\end{quote}

\subsection{ما لم يجده التدقيق}

\textbf{الرياضيات صمدت كاملةً.} أُعيد استنباط القضايا الأربع التي تقوم عليها
الطريقة وأُعيد فحصها: نظرية الشريحة، وتمام تعداد الجبهة، وانفصال التقسيم إلى
\lr{$p$} ابناً، وصحّة الـ\lr{big-$M$}. والأخيرة فُحصت عددياً على \lr{$220$}
زوجاً من النقاط بلا خرقٍ واحد، وبالحدّ مُدرَكاً بالضبط، وهو ما يجعله أصغر
اختيارٍ صحيحٍ لا مجرّد اختيارٍ آمن. فالأحد عشر كلّها كانت في \textbf{إلزام}
الفرضيات، ولا واحدةٌ منها في الاستدلال الذي يستخدمها.

وثلاث ملاحظاتٍ تُغلق الحلقة، وأصوغها ملاحظاتٍ لا افتراضاتٍ لأن منهج التدقيق
هو ألّا يُفترض مثلُ هذا بالضبط:

\begin{enumerate}[leftmargin=1.9em,itemsep=2pt]
\item كلّ رقمٍ منشورٍ أُنتج بتمرير الصندوق الذي يستنبطه الموديل لنفسه، لا
      صندوقاً مكتوباً باليد، فلم يتأثّر أيٌّ منها.
\item مع تشغيل فحص النتيجة الحادية عشرة نجحت المجموعة \lr{$103/103$}، وعملت
      الأمثلة الأربعة التي تستخدم المراجع بلا خطأ. وهذا \textbf{دليلٌ} --- لا
      افتراض --- على أنّ أيّاً من الصناديق الـ\lr{$45$} المكتوبة باليد لم يكن
      يقصّ \D\ قطّ: الفخّ كان في الواجهة البرمجية، لا في أيّ قياسٍ مُعلَن.
\item ودراسة التوسّع ما زالت تُبرهن كل نموذجٍ من \lr{$n=10$} إلى
      \lr{$n=20$} بعد إغلاق النتيجة العاشرة، فالحارس لم يُبطِل الدراسة التي
      تقوم عليه.
\end{enumerate}

\noindent
والحُرّاس دائمون لا وضعُ تنقيحٍ، لأن الكلفة رسمٌ ثابتٌ لكل مسألةٍ لا لكل
مسألةٍ فرعية: الحدودية برنامجٌ صحيحٌ واحد للمسألة كلها؛ والمقامات برنامجٌ لكل
معيارٍ \emph{كسريٍّ فعلاً} ولا شيء للخطّي؛ وصندوق المتغيّرات \lr{$n$} برنامجاً
مُخزَّنةً على المسألة؛ والتحقّق من بذرةٍ هو الفحص الذي كانت البذرة تحتاجه
أصلاً قبل أن تُؤتمَن.

\section{دروس القياس}

\begin{itemize}[leftmargin=1.6em,itemsep=3pt]
\item \textbf{مقياس جودةٍ غير سليم.} قُيِّمت جودة البذرة أولاً بالنسبة
      \lr{$\Phi(s)/\Phi^{*}$}، وهذه \emph{تنقلب إذا كان المقام سالباً}، وهي سالبةٌ في عدّة نماذج: فتُسجّل بذرةٌ أسوأ فوق
      الواحد. والمستخدم الآن هو الفجوة غير السالبة
      \lr{$(\Phi^{*} - \Phi(s))/|\Phi^{*}|$}.
\item \textbf{العدّ الحتمي بدل الثواني.} عدد البرامج الصحيحة وعُقَد التفريع
      والتحديد ومحاوِر \lr{simplex} يُعطي الرقم نفسه على كل تشغيلٍ وكل آلة،
      بخلاف الزمن.
\item \textbf{وسيط النسب لا نسبة المجاميع}، والتعميم من نموذجٍ واحدٍ خطأ ---
      وقعتُ فيه خمس مرات في هذا المشروع. مثالٌ مُصحَّح: نموذجٌ واحدٌ أعطى
      تغطيةً \lr{$35\%$} ونسبةً \lr{$1.00\times$}، وستّة نماذج تُعطي وسيطاً
      \lr{$78\%$} و\lr{$1.35\times$}.
\end{itemize}

\section{الحصيلة}

\begin{enumerate}[leftmargin=1.9em,itemsep=3pt]
\item \textbf{الوصول إلى المجموعة الكفؤة الكاملة \E، لا الجبهة فقط}، بفضل
      نظرية الشريحة ودون دفع فحص كفاءةٍ على نقاط الشرائح. وهذا هو المنتظَر في
      مقترح الرسالة: «مجموعة من الحلول الكفؤة مولَّدة بخوارزميات هجينة».
\item \textbf{هجينٌ مضبوط--ميتاهيريستي يدفع فعلاً}: \lr{$1.36\times$} على
      \lr{$120$} نموذجاً، و\lr{$1.39\times$} على المعايير الكسرية، وكلّ
      العائلات الخمس تدفع --- بأجوبةٍ متطابقة دائماً.
\item \textbf{هجينٌ مضبوط--مضبوط} جُرِّب وقِيس وخَسِر، وهذا مُسجَّلٌ لا مطموس.
\item \textbf{حذف متغيّرات الشريحة}: أكبر مكسبٍ منفرد، \lr{$49\times$} وسيطاً.
\item \textbf{الترتيب هو الآلية}: السطر الواحد الذي يقلب الهجين من رابحٍ إلى
      خاسر، وهو أقبل نتائج المشروع للتعميم.
\item \textbf{صنف دالّة القرار المقبولة}: \E\ الصريحة تقبل \textbf{أيّ
      دالّةٍ قابلةٍ للحساب} --- مجموع نسبٍ، أقصى نسبٍ، صندوقٌ أسود، أو عدّة
      دوالّ دفعةً واحدة --- حيث تشترط طرقُ فضاء القرار صيغةَ النسبة لتبني
      قطوعها. وهذه أقوى ما نملك، لأنها ليست ادّعاءَ سرعة.
\item \textbf{تدقيقٌ أعاد فحص أدوات التحقّق نفسها}، فوجد أحد عشر خللاً، ولم
      يُحرّك رقماً، وأثبت بالملاحظة أنّ النتائج المنشورة سليمة.
\end{enumerate}

\section{ما تشتريه التوليفة: صنف دالّة القرار المقبولة}

كلّ ما سبق آلةٌ للوصول إلى \E. وهذا القسم هو ما \textbf{من أجله} تلك الآلة،
وهو الجزء الذي ندافع عنه أوّلاً --- لأنه لا يقوم على كوننا الأسرع في شيء، بل
هو قولٌ في \textbf{أيّ المسائل يمكن طرحها أصلاً}.

\subsection{المقابلة}

\begin{quote}
\textbf{طريقةُ فضاء القرار تحتاج \emph{بنيةً في} \lr{$\Phi$} ولا تحتاج شيئاً
من \E؛ وهذا الطريق يحتاج \E\ ولا يحتاج شيئاً من \lr{$\Phi$}.}
\end{quote}

\noindent
كلّ طريقةٍ في الأدبيات الكسرية تبني قطوعها وحدودها ومسائلها الفرعية من
\textbf{صيغة النسبة} في \lr{$\Phi$}. وهذا ليس عَرَضاً، بل هو الآلية. فأبدِلْ
\lr{$\Phi$} بـ\textbf{مجموع} نسبتين، فلا تفقد تلك الطريقة كفاءتها --- بل
\textbf{تتوقّف عن الانطباق}، إذ لا نسبةَ واحدةً تُخطَّط أو تُغذّي خطوة
\lr{Dinkelbach}.

\begin{quote}
\textbf{قضية (صنف المقبولات).} لتكن \E\ معلومةً ومُبرهَناً تمامُها. فلأيّ
دالّة \lr{$f$} معرَّفةٍ على \D، تُحسَب \lr{$\max\{f(x) : x \in E(P_D)\}$}
ومجموعةُ معظِّماتها كاملةً \textbf{بالتقييم على \E}. ولا يُشترط خطّيةٌ ولا
اطّرادٌ ولا اتصالٌ ولا اشتقاقٌ ولا صيغةُ نسبة.
\end{quote}

\noindent
والبرهان هو \textbf{منتهاهيةُ \E} ولا شيء غيرها. ونذكرها مع ذلك لأن المضمون
ليس في القضية، بل في \textbf{أنّنا دفعنا ثمن \E} بسعرٍ يجعل الدفع مُجدياً ---
والقضية هي ما يشتريه ذلك السعر. وهذا أيضاً سببُ أنّ نتائج الشريحة والهجين
تُهمّ أبعدَ من هوامشها: هي ما يجعل فرضيةَ هذه القضية \textbf{في المتناول}.

\begin{table}[h]
\centering
\caption{ستّ صيغٍ لدالّة القرار، كلٌّ منها مفحوصٌ مقابل تعظيمٍ غاشمٍ مستقل.}
\begin{tabular}{lcr}
\toprule
\textarabic{دالّة القرار} & \textarabic{تقبلها طرق فضاء القرار؟}
  & \textarabic{القيمة المثلى} \\
\midrule
\lr{$\Phi_1$}              & \textarabic{نعم --- الحالة الكلاسيكية} & \lr{$5/17$} \\
\lr{$\Phi_1 + \Phi_2$}     & \textarabic{لا} & \lr{$47/35$} \\
\lr{$\max(\Phi_1,\Phi_2)$} & \textarabic{لا} & \lr{$6/5$} \\
\lr{$\Phi_1 \cdot \Phi_2$} & \textarabic{لا} & \lr{$6/35$} \\
\lr{$(x_1^2+1)/(x_2^2+1)$} & \textarabic{لا} & \lr{$10$} \\
\textarabic{مجزَّأة / صندوقٌ أسود} & \textarabic{لا} & \lr{$3/11$} \\
\bottomrule
\end{tabular}
\end{table}

\subsection{عدّة دوالّ قرارٍ دفعةً واحدة}

تعظيم \textbf{دالّتين} خطّيتين كسريتين فوق مجموعةٍ كفؤةٍ صحيحة مسألةٌ قائمةٌ
بذاتها في تلك الأدبيات، لها تفريعُها وتحديدُها وقطوعُها. وهنا هي
\textbf{مُرشِّحُ هيمنةٍ} فوق مجموعةٍ في اليد، والكلفة لا تتعلّق بعدد الدوالّ
المطلوبة ولا بماهيّتها: \lr{$m$} نسبةً، أو نسبةً وتربيعية، أو عشرة صناديق
سوداء --- المسحُ نفسه. وعلى النموذج المنشور يُبقي مُرشِّحُ النسبتين \lr{$4$}
من \lr{$7$} نقاطٍ كفؤة، وهو بالضبط جواب القوة الغاشمة.

\subsection{نقطة التعادل، ومتى يكون هذا الخيار خاطئاً}

مقايضةٌ لا وجبةٌ مجانية، فتُقاس. والبديل الأمين هو تشغيل الطريقة المضبوطة
\textbf{مرّةً لكلّ دالّة قرار}.

\begin{table}[h]
\centering
\caption{كم دالّةَ قرارٍ قبل أن يُجدي بناءُ \E؟ بعدد \textbf{البرامج
الفرعية}، ثلاثون نموذجاً لكل تهيئة. الجدول مُولَّدٌ من ملف بيانات الحملة لا
منسوخاً باليد.}
\begin{tabular}{lcccl}
\toprule
\textarabic{التهيئة} & \textarabic{نماذج} & \textarabic{مُقصى}
  & \lr{$|E(P_D)|$} & \textarabic{نقطة التعادل} \\
\midrule
\lr{$n=5$, $p=3$} & \lr{$30$} & \lr{$0$} & \lr{$4$} & \lr{$\mathbf{1.8}$ $[1.6,\ 2.5]$} \\
\lr{$n=6$, $p=3$} & \lr{$30$} & \lr{$0$} & \lr{$6$} & \lr{$\mathbf{2.2}$ $[2.0,\ 2.5]$} \\
\lr{$n=6$, $p=4$} & \lr{$27$} & \lr{$3$} & \lr{$9$} & \lr{$\mathbf{3.1}$ $[2.1,\ 4.2]$} \\
\lr{$n=7$, $p=3$} & \lr{$30$} & \lr{$0$} & \lr{$8$} & \lr{$\mathbf{2.6}$ $[2.0,\ 3.3]$} \\
\midrule
\textarabic{\textbf{الكل}} & \lr{$117$} & & & \lr{$\mathbf{2.4}$ $[1.8,\ 2.9]$} \\
\bottomrule
\end{tabular}
\end{table}

\begin{quote}
\textbf{هذا الجدول ليس المقارنة بين الطريقتين.} إنه يُسعّر \textbf{منتَجاً
مقابل منتَجٍ آخر} --- بناءَ \E\ كاملةً مقابل الإجابة عن سؤالٍ واحد --- والأول
أغلى بطبيعة الحال. أمّا المقارنة وجهاً لوجه، والطريقتان تُجيبان عن
\lr{$(P_E)$} نفسها، فهي الجدول \lr{1}: \textbf{فضاء المعايير يفوز في كل
صفّ}، من \lr{$1.52\times$} عند \lr{$p=2$} إلى \lr{$14.57\times$} عند
\lr{$p=6$}، باطّرادٍ مع \lr{$p$}، والاقتطاع في صالح الطريقة الأخرى، والقيمة
المثلى متطابقة دائماً. وما يلي يُسعّر \textbf{المنتَج الزائد}، لا عجزاً في
المشترَك.
\end{quote}

\noindent
الوسيط المجمَّع \lr{$\mathbf{2.4}$}، والمدى الربيعي \lr{$[1.8,\ 2.9]$} على
\lr{$117$} نموذجاً، وطرفاه \lr{$1.2$} و\lr{$7.2$}.

\textbf{والإحصاءة بعدد البرامج الفرعية لا بالثواني، وهذا ما يسمح بذكرها بهذا
الضيق.} فالمسح فوق \E\ صريحة يحلّ \textbf{صفر} برنامجٍ صحيح، فالنسبة هي
حرفياً «برامج البناء ÷ برامج تشغيلٍ مضبوطٍ واحد»: كمّيةٌ \textbf{حتميّة}،
واحدةٌ على كل آلةٍ وكل إعادة، لا تمسّها أرضية الضوضاء \lr{$\pm 30\%$} التي لا
يُصلحها أيّ عدد من التكرارات. وهذا أوضح مثالٍ في العمل كلّه على درس القياس:
عُدَّ العملَ الحتميّ لا الثواني --- إذ حوّل رقماً يمتدّ ثلاث رُتبٍ عشرية إلى
رقمٍ يمتدّ بعامل ستة.

\paragraph{ورقمٌ مسحوب.}
كانت نسخةٌ سابقة من هذا القسم تذكر وسيطاً \lr{$13$} بمدًى من \lr{$1$} إلى
\lr{$2117$}. \textbf{وهو مسحوب.} جاء من ثمانية نماذج صنعتُها يدوياً ببنية
معايير متعارضةٍ عمداً و\lr{$|E(P_D)|$} بين \lr{$17$} و\lr{$105$}؛ بينما
عائلة الحملة \lr{$|E(P_D)|$} فيها بين \lr{$4$} و\lr{$9$}، والمجموعة الصغيرة
رخيصةُ البناء. ولم يكن الرقم القديم قياساً فاسداً --- هو ما فعلته تلك النماذج
الثمانية --- لكنه ذُكِر بوصفه \textbf{نقطة التعادل} وهو نقطة تعادل
\textbf{عائلةٍ واحدة}. والتعميم من عائلةٍ واحدة هو الخطأ الذي يسجّل هذا العمل
وقوعه فيه خمس مرات، وهذه سادستها --- كشفها إعادةُ القياس بمعيار الحملة، لا
الانتباه.

وأُقصيت ثلاثة أشواطٍ من \lr{$120$}، كلّها عند \lr{$p=4$}، حين استنفد تشغيلٌ
مضبوطٌ ميزانيته. وهذا الإقصاء \textbf{في غير صالح} ما ندافع عنه هنا: التشغيل
المضبوط المتجاوز هو تشغيلٌ غالٍ، والغالي هو بالضبط ما يُجدي معه بناءُ \E\
أسرعَ ما يكون. فـ\lr{$2.4$} حدٌّ \textbf{أعلى} للحقيقة لا أدنى.

\paragraph{النصف غير المُجامِل، ويُذكَر أولاً.}
لدالّة قرارٍ \textbf{واحدةٍ} خطّيةٍ كسرية، تكون طريقةُ فضاء القرار غالباً
\textbf{الصفقةَ الأفضل}، وصفّان أعلاه دون التعادل عند دالّةٍ واحدة. فهي لا
تبني \E\ فلا يمسّها حجمُها؛ وهذا الطريق يبنيها فيتقيّد بها. والطريقتان
تفشلان على نماذج مختلفة --- على \lr{$f$} بلا بنية، وعلى \E\ أكبر من أن
تُعَدّ --- ولهذا تُعرَضان \textbf{جنباً إلى جنب لا مُرتَّبتين}.

\paragraph{والنصف الذي لا يتعلّق بالقياس أصلاً.}
لأيّ \lr{$f$} خارج الصنف الخطّي الكسري، نقطةُ التعادل \lr{$1$}، إذ
\textbf{لا منافسَ يُستهلَك ثمنُه}. وهذا ليس ادّعاءَ سرعة، ولا يَنقُضه أيُّ
توقيت.

\section{موقع العمل من الأدبيات}

جزءان من هذا البناء \textbf{ليسا جديدين، ولا نَدّعيهما}:

\begin{itemize}[leftmargin=1.6em,itemsep=2pt]
\item \textbf{تفكيك فضاء المعايير إلى صناديق}: راسخٌ عند
      \lr{Boland, Charkhgard \& Savelsbergh} في \lr{EJOR} سنة \lr{2017}،
      يعملون في فضاء المعايير كامل البُعد ويفكّكونه في كل دورة --- وهو بنيوياً
      طريقتنا بهدفٍ ثانويٍّ \textbf{خطّي}. وطرق \lr{balanced box}
      و\lr{boxed line} و\lr{triangle splitting} هي سوابقها ثنائية الهدف.
\item \textbf{نمط التهجين ميتاهيريستي--مضبوط}: له تصنيفٌ معروف، و«شغّل
      ميتاهيريستي لتقريب جبهة \lr{Pareto} ثم دع الطريقة المضبوطة تُنقّحها» هو
      بالضبط مخطّطنا.
\end{itemize}

\noindent
وما لا يبدو أنّ الأدبيات الكسرية تفعله هو \textbf{العمل في فضاء المعايير}.
فتعظيم دالّةٍ خطّيةٍ كسريةٍ فوق مجموعةٍ كفؤةٍ صحيحة له أدبياتٌ مستمرّة ---
\lr{Drici, Ouail \& Moula\"i} في \lr{Annals of OR 2018}، وورقة
\lr{RAIRO-OR 2015}، و\lr{Ouail \& Chergui}، و\lr{Chaiblaine, Moula\"i \&
Cherfaoui}، وورقة \lr{Optimization 2021} للأهداف التربيعية، وأعمال
\lr{Younsi-Abbaci \& Moula\"i} العشوائية والفترية \lr{(2021--2023)} ---
و\textbf{كلّها تبحث في فضاء القرار} بالتفريع والتحديد أو القطوع.

فالذي يبقى غيرَ مطروقٍ على ما يبدو \textbf{أضيقُ من الطريقة}: هو
\textbf{التوليفة} نفسها، ومعها ما يجعل النقل مضبوطاً لا تقريبياً --- الصيغة
\lr{$e_k$} الصحيحة القيمة التي تُبدِل تقديرَ أصغر خطوةٍ بـ\lr{$e_k \geqslant 1$}
بالضبط، ونتائج الشريحة.

\paragraph{وحَدّان يجب ذكرهما.}
\textbf{لم يُقرأ أيُّ نصٍّ كاملٍ من هذه القائمة}: جُمِعت من المستخلصات
والفهارس، والبيئة التي أُنجز فيها العمل لا تصل إلى \lr{arxiv.org} ولا
\lr{springer} ولا \lr{sciencedirect}. فطريقةٌ يسمّيها مستخلصُها «تفريعاً
وتحديداً» قد تفعل في فضاء المعايير ما لا يذكره المستخلص. \textbf{وورقةٌ واحدة
تمسّ نتيجة المجموعة الكاملة مباشرةً ويجب قراءتها قبل وصفها بالجِدّة}:
\lr{``Generating the efficient set of multi-objective integer linear plus
linear fractional programming problems''}، \lr{Annals of Operations Research}
مجلد \lr{\textbf{296}} \lr{(2021)}، \lr{pp.~735--753}.

\textbf{وهذا بيانٌ عمّا وجدنا وما لم نجد، لا ادّعاءُ جِدّة.}

\section{ما بقي}

\begin{itemize}[leftmargin=1.6em,itemsep=2pt]
\item ملء \lr{\texttt{\textbackslash author\{\}}} في الوثائق الخمس، واسم صاحب
      الرخصة.
\item قائمة المراجع في المقالة: لم أختلقها، فهي عليك.
\item قراءة النصوص الكاملة للأوراق المذكورة في القسم السابق، وأوّلها ورقة
      \lr{Annals of OR 296 (2021)} وورقة \lr{EJOR 2017} --- فهما أخطر
      ورقتين على أيّ ادّعاءٍ بالجِدّة، ولا يمكن الوصول إليهما من هنا.
\item دمج طلب الدمج وحذف الفرع (ممنوعٌ من جهة الوكيل، خطأ \lr{403}).
\end{itemize}

\section{الوثائق التفصيلية}

\begin{table}[h]
\centering
\caption{الوثائق بالإنجليزية في مجلّد \lr{\texttt{docs/}}.}
\begin{tabular}{lcl}
\toprule
\textarabic{الملف} & \textarabic{صفحات} & \textarabic{المحتوى} \\
\midrule
\lr{\texttt{hybrid-article.tex}}        & \lr{$19$} & \textarabic{المقالة المستقلّة} \\
\lr{\texttt{hybrid-methods-note.tex}}   & \lr{$31$} & \textarabic{المذكرة التقنية الكاملة} \\
\lr{\texttt{results.tex}}               & \lr{$17$} & \textarabic{النتائج والرسوم} \\
\lr{\texttt{algorithm-walkthrough.tex}} & \lr{$7$}  & \textarabic{الخوارزمية سطراً سطراً} \\
\lr{\texttt{worked-examples.tex}}       & \lr{$6$}  & \textarabic{مثالان بمجهولين وبثلاثة} \\
\bottomrule
\end{tabular}
\end{table}

\noindent
والتنفيذ \lr{Python} خالص، بلا أيّ اعتمادٍ خارجي، بحسابٍ كسريٍّ مضبوطٍ في كل
موضع ولا فاصلةٍ عائمةٍ في أيّ موضع، مع \lr{$106$} اختباراً تُشغَّل على كل
إيداع.

\end{document}
